\begin{figure}[H]
\centering
\begin{tikzpicture}[>=Latex,scale=0.92]
    \coordinate (A) at (-4.3,0);
    \coordinate (S) at (-6.0,0);
    \coordinate (D) at (-4.3,-2.15);
    \coordinate (LT) at (-3.0,1.15);
    \coordinate (RT) at (4.5,1.15);
    \coordinate (LB) at (-3.0,-1.15);
    \coordinate (RB) at (4.5,-1.15);

    % Fuente, divisor y detector
    \draw[black,very thick,fill=gray!20] (S) circle (0.3);
    \draw[black,thick] ($(S)+(-0.16,-0.16)$) -- ($(S)+(0.16,0.16)$);
    \draw[black,thick] ($(S)+(-0.16,0.16)$) -- ($(S)+(0.16,-0.16)$);
    \node[above=5pt] at (S) {fuente};

    \draw[blue!75!black,line width=3pt]
        ($(A)+(-0.38,-0.38)$) -- ($(A)+(0.38,0.38)$);
    \node[above left=4pt] at (A) {$A$};
    \node[font=\small,align=center] at (-5.15,-0.85)
        {divisor\\semitransparente};

    \draw[black,very thick,fill=gray!15]
        ($(D)+(-0.52,-0.24)$) rectangle ($(D)+(0.52,0.24)$);
    \foreach \x in {-0.32,-0.1,0.12,0.34}
        \draw[black,thin] ($(D)+(\x,-0.21)$) -- ($(D)+(\x,0.21)$);
    \node[below=5pt] at (D) {detector};

    % Tubos con agua
    \draw[black,very thick,fill=blue!8]
        ($(LT)+(0,-0.34)$) rectangle ($(RT)+(0,0.34)$);
    \draw[black,very thick,fill=blue!8]
        ($(LB)+(0,-0.34)$) rectangle ($(RB)+(0,0.34)$);
    \node[blue!75!black,font=\small,fill=white,inner sep=1pt]
        at (0.75,1.42) {agua};
    \node[blue!75!black,font=\small,fill=white,inner sep=1pt]
        at (0.75,-1.42) {agua};

    % Flujo opuesto en los tubos
    \foreach \x in {-2.3,-0.8,0.7,2.2,3.7}
        \draw[->,blue!75!black,thick]
            (\x,1.15) -- ++(0.75,0);
    \foreach \x in {-1.55,-0.05,1.45,2.95,4.45}
        \draw[->,blue!75!black,thick]
            (\x,-1.15) -- ++(-0.75,0);
    \node[blue!75!black,font=\small] at (4.05,1.7) {$\vec v$};
    \node[blue!75!black,font=\small] at (-2.55,-1.7) {$-\vec v$};

    % Espejos de retorno
    \draw[black,line width=3.5pt]
        ($(RT)+(0.18,-0.32)$) -- ($(RT)+(0.62,0.32)$);
    \draw[black,line width=3.5pt]
        ($(RB)+(0.18,0.32)$) -- ($(RB)+(0.62,-0.32)$);
    \node[font=\small,align=center] at (5.85,0)
        {espejos de\\retorno};

    % Haz incidente y salida al detector
    \draw[->,red!75!black,very thick]
        ($(S)+(0.32,0)$) -- ($(A)+(-0.18,0)$);
    \draw[->,red!75!black,very thick]
        ($(A)+(0,-0.18)$) -- ($(D)+(0,0.26)$);

    % Camino a favor del flujo (línea continua)
    \draw[red!75!black,very thick]
        ($(A)+(0.18,0.12)$) -- ($(LT)+(0,-0.12)$);
    \draw[->,red!75!black,very thick]
        ($(LT)+(0,-0.12)$) -- ($(RT)+(0,-0.12)$);
    \draw[red!75!black,very thick]
        ($(RT)+(0.38,-0.12)$) -- ($(RB)+(0.38,0.12)$);
    \draw[->,red!75!black,very thick]
        ($(RB)+(0,0.12)$) -- ($(LB)+(0,0.12)$);
    \draw[red!75!black,very thick]
        ($(LB)+(0,0.12)$) -- ($(A)+(0.18,-0.12)$);

    % Camino en contra del flujo (línea discontinua)
    \draw[red!75!black,thick,dashed]
        ($(A)+(0.08,-0.2)$) -- ($(LB)+(0,-0.12)$);
    \draw[->,red!75!black,thick,dashed]
        ($(LB)+(0,-0.12)$) -- ($(RB)+(0,-0.12)$);
    \draw[red!75!black,thick,dashed]
        ($(RB)+(0.18,-0.12)$) -- ($(RT)+(0.18,0.12)$);
    \draw[->,red!75!black,thick,dashed]
        ($(RT)+(0,0.12)$) -- ($(LT)+(0,0.12)$);
    \draw[red!75!black,thick,dashed]
        ($(LT)+(0,0.12)$) -- ($(A)+(0.08,0.2)$);

    \node[red!75!black,font=\small] at (0.8,2.05)
        {haz a favor del flujo};
    \node[red!75!black,font=\small] at (0.8,-2.05)
        {haz en contra del flujo};

    % Longitud útil
    \draw[gray!65,densely dotted] (LT) -- (-3,2.55);
    \draw[gray!65,densely dotted] (RT) -- (4.5,2.55);
    \draw[<->,gray!75!black,thick]
        (-3,2.4) -- (4.5,2.4)
        node[midway,fill=white,inner sep=2pt] {$L$};
\end{tikzpicture}
\caption{Esquema del experimento de Fizeau de 1851. El agua circula en sentidos opuestos por dos tubos paralelos. Los dos haces recorren el mismo circuito óptico en direcciones contrarias: uno se propaga a favor del flujo y el otro en contra, produciendo una diferencia de fase al recombinarse.}
\label{fig:fizeau1851}
\end{figure}
